\documentclass[english,handout]{mlutalk}

\title{Stance Classification for Key-Point Analysis}
% \title{%
%   Modern Talking: Key-Point Analysis \\
%   using Modern Natural Language Processing
% }
\subtitle{Natural Language Processing, Summer Semester 2021}
\author{Max Henze \and Hanh Luu \and Jan Heinrich Reimer}
\institute{Martin Luther University Halle-Wittenberg}
\date{\today}
\titlegraphic{\includegraphics[width=3cm]{figures/mlu-halle}}

\addbibresource{../literature/literature.bib}

\usepackage{tikz}
\usepackage{listings}
\usepackage{xspace}
\usepackage{biblatex}

\newcommand{\Bert}{\textsc{Bert}\xspace}
\newcommand{\ArgKP}{\mbox{ArgKP}\xspace}
\newcommand{\ArgQ}{\mbox{IBM-ArgQ-Rank-30kArgs}\xspace}

\begin{document}

\titleframe

\begin{frame}{Stance Classification}
  TODO: motivation and problem description.
\end{frame}

\begin{frame}{Approaches}
  TODO: general framework and basic approaches.
\end{frame}

\begin{frame}{Fake News Challenge Stance Detection~\cite{HanselowskiSSCC2018}}
  TODO
\end{frame}

\begin{frame}{Argument Stance Classification Using \Bert~\cite{OllingerDSBS2020}}
  TODO
\end{frame}

\begin{frame}{Cross-topic Argument Mining~\cite{StabMSRG2018}}
  TODO
\end{frame}

\begin{frame}{Comparison}
  TODO: compare approaches
\end{frame}

\begin{frame}{Stance in Key Point Matching}
  TODO: how can we use stance classification for matching key points?
\end{frame}

\begin{frame}{Conclusion}
  TODO: what to do next?
\end{frame}

\begin{frame}{Future Work}
  TODO: do we need future work?
  \thankyou
\end{frame}

\appendix
\section{\appendixname}

\bibliographyframe

\end{document}
